\documentclass[11pt,a4paper]{article}

% ---------- Packages ----------
\usepackage[margin=1in]{geometry}
\usepackage{times}
\usepackage[T1]{fontenc}
\usepackage{titlesec}
\usepackage{fancyhdr}
\usepackage{booktabs}
\usepackage{array}
\usepackage{colortbl}
\usepackage{longtable}
\usepackage{graphicx}
\usepackage{caption}
\usepackage{enumitem}
\usepackage{xcolor}
\usepackage{hyperref}
\usepackage{amsmath}
\usepackage{amssymb}
\usepackage{tikz}
\usepackage{tcolorbox}
\usepackage{tabularx}
\usepackage[super,comma,sort&compress]{natbib}
\tcbuselibrary{skins,breakable}

% ---------- Colors ----------
\definecolor{havennavy}{HTML}{1F3864}
\definecolor{havensteel}{HTML}{2E5395}
\definecolor{havenlight}{HTML}{DCE6F1}
\definecolor{havenmute}{HTML}{595959}

% ---------- Section styling ----------
\titleformat{\section}
  {\normalfont\Large\bfseries\color{havennavy}}
  {\thesection}{0.6em}{}[{\color{havensteel}\titlerule[0.9pt]}]
\titlespacing{\section}{0pt}{18pt}{10pt}

\titleformat{\subsection}
  {\normalfont\large\bfseries\color{havensteel}}
  {\thesubsection}{0.6em}{}
\titlespacing{\subsection}{0pt}{14pt}{6pt}

\titleformat{\subsubsection}
  {\normalfont\normalsize\bfseries\itshape\color{havennavy}}
  {\thesubsubsection}{0.6em}{}
\titlespacing{\subsubsection}{0pt}{10pt}{4pt}

% ---------- Header / Footer ----------
\pagestyle{fancy}
\fancyhf{}
\renewcommand{\headrulewidth}{0.4pt}
\fancyhead[R]{\small\itshape\color{havenmute} HAVEN Project Documentation}
\fancyfoot[C]{\small\color{havenmute} Page \thepage\ of \pageref{LastPage}}
\usepackage{lastpage}

% ---------- Callout box ----------
\newtcolorbox{notebox}[1]{
  breakable, enhanced,
  colback=havenlight, colframe=havensteel,
  boxrule=0.5pt, leftrule=3pt,
  fonttitle=\bfseries\color{havennavy},
  title={#1},
  arc=1mm, top=4pt, bottom=4pt, left=8pt, right=8pt
}

% ---------- Misc ----------
\setlist[itemize]{leftmargin=1.4em, topsep=2pt, itemsep=3pt}
\setlist[enumerate]{leftmargin=1.6em, topsep=2pt, itemsep=3pt}
\renewcommand{\arraystretch}{1.3}
\hypersetup{colorlinks=true, linkcolor=havensteel, urlcolor=havensteel, citecolor=havensteel, pdftitle={HAVEN Project Documentation}}

\newcommand{\citeref}[1]{}

\begin{document}

% ==================== TITLE PAGE ====================
\begin{titlepage}
\centering
\vspace*{3cm}
{\fontsize{48}{54}\selectfont\bfseries\color{havennavy} HAVEN\par}
\vspace{0.4cm}
{\Large\bfseries\color{havensteel} Hostile-Aware Verifiable Election Network\par}
\vspace{0.2cm}
{\large\itshape\color{havenmute} Comprehensive Technical Documentation\par}
\vspace{1.2cm}
{\color{havensteel}\rule{0.6\textwidth}{0.8pt}}
\vfill
{\normalsize\color{havenmute}
A cryptographically secure leader-election and communication protocol\\
for decentralized wireless sensor networks in hostile environments\par}
\vspace{1.5cm}
{\small\color{havenmute} \today \par}
\end{titlepage}

% ==================== TOC ====================
\tableofcontents
\thispagestyle{fancy}
\newpage

% ==================== 1. INTRODUCTION ====================
\section{Introduction}

The HAVEN (Hostile-Aware Verifiable Election Network) project is a research and implementation initiative aimed at designing a cryptographically secure leader-election and communication protocol.\citeref{1,4} It is specifically tailored for decentralized Wireless Sensor Networks (WSNs) operating in hostile environments such as battlefields, disaster-recovery zones, and remote industrial asset-monitoring sites, where there is no base station or trusted infrastructure.\citeref{1,4}

\subsection{The Core Challenge}

The primary challenge addressed by HAVEN is achieving authenticated, tamper-resistant leader election without relying on trusted infrastructure, public-key cryptography, or pre-shared pairwise keys.\citeref{1,4} The protocol must use only symmetric primitives on severely resource-constrained embedded nodes.\citeref{4}

\subsection{Threat Model}

The network operates under an active adversary model.\citeref{1} The adversary is capable of eavesdropping, packet injection and dropping, replay attacks, Man-in-the-Middle (MITM) attacks, and fake-node insertion.\citeref{1} However, the nodes themselves do not act Byzantine.\citeref{1}

\subsection{Limitations of Existing Protocols}

Traditional consensus and leader-election models are impractical for this environment due to a range of limitations:\citeref{1}

\begin{longtable}{@{}>{\bfseries}p{3.2cm}p{11cm}@{}}
\toprule
\rowcolor{havennavy}
\textcolor{white}{Approach} & \textcolor{white}{Limitation} \\
\midrule
\endhead
Paxos & Too complicated for implementation on weak hardware; the total number of nodes is not pre-known. \\
Raft & Imposes heavy overhead, unsuitable for resource-constrained devices. \\
Ring & Very fragile and susceptible to node drops. \\
Bully & Power-intensive and requires a dense network. \\
Traditional key exchange & Relies on heavy PKI or Diffie--Hellman; a standard DH exchange can take nearly a minute on a typical IoT device, making it vulnerable to Denial-of-Service (DoS) attacks. \\
\bottomrule
\end{longtable}

\subsection{Solution Objective}

The objective is to elect a verifiable leader using lightweight symmetric cryptographic primitives on resource-constrained embedded nodes.\citeref{1}

\begin{notebox}{Protocol Convention}
Unless a field or packet is explicitly described elsewhere in this document as encrypted, it is assumed to be sent in the clear and accompanied by an HMAC over that field, generated with the appropriate derived key for the context (e.g.\ $K_{master}$, $K_{node}$, or $K_{AB}$). This default holds for every packet type in Sections~3 and beyond, including any not yet listed with an explicit HMAC field.
\end{notebox}

% ==================== 2. KEY PRE-DISTRIBUTION ====================
\section{Key Pre-distribution and TurboBlom}

To establish secure communication without the overhead of real-time mathematical key negotiation, constrained networks utilize Key Pre-distribution Schemes (KPS).\citeref{4} The HAVEN project implements a modified version of the TurboBlom scheme.

\subsection{The Classic Blom Scheme}

The classic Blom scheme provides 100\% connectivity and optimal resilience against node capture up to a threshold $\lambda$.\citeref{4} It uses finite-field algebra, where a trusted server generates a public generator matrix $G$ and a secret symmetric matrix $D$.\citeref{4}

The share matrix $A$ is computed as
\begin{equation*}
A = (D \cdot G)^{T}.
\end{equation*}
\citeref{4} Nodes are preloaded with a row from $A$ and a column from $G$, and calculate a shared key using an inner product.\citeref{4} The major drawback is the high number of computationally expensive field multiplications required.\citeref{4}

\subsection{The Academic TurboBlom Scheme}

The academic TurboBlom scheme (Khabbazian et al., 2024) eliminates field multiplications by restricting the public generator matrix $G$ to a confined subset, such as $\{0,1\}$ or powers of 2 (when using Mersenne primes).\citeref{4} This converts heavy multiplications into fast additions or bit-shifts, while maintaining strict $\lambda$-collusion resistance.\citeref{4}

\subsection{Project HAVEN's Implementation of ``TurboBlom''}

The implementation in the HAVEN MVP deviates from the academic TurboBlom scheme. It is a probabilistic, combinatorial random key pre-distribution scheme optimized with bitwise operations and hashing, conceptually similar to the Eschenauer--Gligor scheme.\citeref{4}

\subsubsection*{Architectural Shift: From Storage-Bound to Compute-Bound}

The architecture evolved from a storage-bound model (storing a matrix of raw bytes) to a compute-bound Key Derivation Function (KDF) using HMAC-SHA256.\citeref{4}

\begin{itemize}
\item \textbf{Master Key} --- every node holds a single 32-byte, network-wide $K_{master}$ and a list of assigned integer `spaces'.
\item \textbf{Spaces} --- a `space' is a virtual cryptographic domain identifier (e.g., an integer up to 65{,}535) that acts as a variable input seed to the HMAC algorithm.
\end{itemize}

\subsubsection*{Key Derivation Process}

\begin{enumerate}
\item \textbf{Intersection} --- nodes exchange their assigned space lists and find the overlapping indices: $S_{shared} = S_A \cap S_B$. If there is no overlap, they cannot establish a direct key. The intersection is sorted to ensure deterministic execution.
\item \textbf{Symmetry Enforcement} --- to ensure both nodes generate identical HMAC inputs, node IDs are sorted: $id_{low} = \min(node\_id_a, node\_id_b)$ and $id_{high} = \max(node\_id_a, node\_id_b)$.
\item \textbf{HMAC Generation Loop} --- for each shared space index, a 6-byte message is constructed: $Message = id_{low} \,\|\, id_{high} \,\|\, space\_idx$. This message and $K_{master}$ are fed into HMAC-SHA256; the first 16 bytes of the digest form a sub-key.
\item \textbf{XOR Accumulation} --- an accumulator is initialized with 16 zero-bytes. Each generated sub-key is sequentially XORed into this accumulator. The final result is the pairwise key $K_{AB}$.
\end{enumerate}

\subsubsection*{Trade-offs and Vulnerabilities}

This approach reduces matrix storage to zero and allows infinite scalability.\citeref{4} However, it introduces a fatal flaw regarding Node Capture Attacks: because every node holds the same $K_{master}$, capturing a single node instantly compromises the entire network, allowing the attacker to eavesdrop, impersonate, and decrypt historical data.\citeref{4}

\begin{notebox}{Mitigation}
To mitigate this in a real-world scenario, $K_{master}$ must be stored in a Hardware Security Module (HSM) or Secure Enclave, which prevents the software from reading the key directly.
\end{notebox}

% ==================== 3. NETWORK PROTOCOL ====================
\section{Network Protocol and Communication}

\subsection{Modified Gossip Election}

HAVEN employs a modified Gossip Election protocol for system architecture.\citeref{1} A standard gossip protocol relies on random peer selection to spread information.\citeref{1} Since HAVEN uses a Chordal Ring topology, peer selection is limited, removing the need for random selection.\citeref{1}

The protocol propagates candidate priority across the mesh, suppresses redundant broadcasts, and locks votes on a timeout, ensuring convergence without a central coordinator.\citeref{1}

\begin{figure}[htbp]
\centering
\includegraphics[width=0.92\textwidth]{figs/fig_chordal_ring.png}
\caption{Two independently elected segments of a Chordal Ring topology mid-search. Each ring shows its own leader (green/amber marker) and a short run of in-flight search packets (cyan) propagating along the ring toward it.}
\label{fig:chordal-ring}
\end{figure}

Figure~\ref{fig:chordal-ring} shows the Chordal Ring structure that motivates this design choice: because each node's neighbours are fixed by its position on the ring (rather than chosen at random), search and election traffic only ever needs to travel along a small, predictable set of links to reach a leader.

\subsection{Leader Election and Consensus}

\begin{enumerate}
\item \textbf{Self-Vote Tracking} --- a self-voting candidate tracks the echo instances of its election packet across adjacent nodes.
\item \textbf{Threshold Validation} --- when the echo count reaches a predefined threshold (fine-tuned via simulation), a safety timeout is initiated.
\item \textbf{Claimant Broadcast} --- after the timeout, the candidate broadcasts a Claimant Packet.
\item \textbf{Consensus Finalization} --- if no objections occur during the claim window, the node transitions to the Leader state.
\end{enumerate}

\begin{figure}[htbp]
\centering
\includegraphics[width=0.92\textwidth]{figs/fig_leader_active.png}
\caption{Simulation snapshot of the leader-election protocol in the \texttt{LEADER\_ACTIVE} phase. Node N8 (highlighted in green) holds the leader role, shown with its battery level, composite score, and term progress relative to the configured epoch limit.}
\label{fig:leader-active}
\end{figure}

Figure~\ref{fig:leader-active} shows a live simulation of the protocol after convergence: node scores (a composite of battery level and node degree) determine the elected leader, whose term is tracked in epochs and re-evaluated as it progresses.

\subsection{Session Caching and Re-authentication}

To avoid the computationally heavy HMAC HELLO handshake for every message, nodes implement a stateful architecture using a Neighbor Table.\citeref{4}

\begin{itemize}
\item \textbf{Neighbor Profile} --- stores the Authentication Status, a Countdown Timer (measured in `epochs'), a Penalty Tracker (miss count), and the cached derived key ($K_{AB}$).
\item \textbf{Session Expiry} --- a background heartbeat reduces the countdown timer each epoch; when the timer hits 0, the session is invalidated.
\item \textbf{Re-authentication} --- a new HELLO handshake is required. If successful, the timer is reset and penalties are cleared; if it fails, the neighbor's profile and cached keys are completely deleted.
\end{itemize}

\subsection{Link Failure and Network Partitioning}

\begin{itemize}
\item \textbf{Disconnection Detection} --- severed links trigger a `Find Leader' packet containing the node's ID and the lost neighbor's ID.
\item \textbf{Loop Termination} --- if a node receives the search packet from both sides of the severed link, the search terminates, confirming the nodes are in the same segment.
\item \textbf{Segmentation Detection} --- if one node finds the leader but the other times out, a network partition is confirmed.
\item \textbf{Isolated Re-Election} --- the disconnected segment triggers a new local leader election.
\end{itemize}

\begin{figure}[htbp]
\centering
\includegraphics[width=0.92\textwidth]{figs/fig_search_propagation.png}
\caption{Simulation snapshot of the \texttt{SEARCH\_PROPAGATING} phase across a randomly generated mesh topology. Cyan edges denote in-flight search (`Find Leader') packets; the green node marks a leader acknowledgment already received by that segment.}
\label{fig:search-propagation}
\end{figure}

Figure~\ref{fig:search-propagation} illustrates the flood-search mechanism in progress: search packets (cyan) radiate outward from the point of disconnection while a node that has already located the leader (green) anchors that segment of the network.

\subsection{Network Merging and Leader Resignation}

\begin{itemize}
\item \textbf{Join Request Delegation} --- nodes receiving external join requests forward them to the active leader.
\item \textbf{Resignation Threshold} --- if join requests exceed a limit, the active leader resigns to merge the networks.
\item \textbf{Global Re-Election} --- a unified election begins across the new combined topology.
\end{itemize}

\subsection{Packet Structures}

\subsubsection*{Election Packet (40 Bytes)}

\begin{longtable}{@{}>{\bfseries}p{3.4cm}p{2cm}p{8.5cm}@{}}
\toprule
\rowcolor{havennavy}
\textcolor{white}{Field} & \textcolor{white}{Size} & \textcolor{white}{Description} \\
\midrule
\endhead
Score & 4 Bytes & Score metric for better leader choice. \\
Origin ID & 2 Bytes & ID of the node owning the score. \\
Sender ID & 2 Bytes & ID of the neighbour forwarding the packet. \\
Score HMAC & 16 Bytes & HMAC of Score generated by the origin node. \\
Origin ID HMAC & 16 Bytes & HMAC of Origin ID generated by the origin node. \\
\bottomrule
\end{longtable}

\subsubsection*{Leader Heartbeat (38 Bytes + 2 Bytes Padding)}

\begin{longtable}{@{}>{\bfseries}p{3.4cm}p{2cm}p{8.5cm}@{}}
\toprule
\rowcolor{havennavy}
\textcolor{white}{Field} & \textcolor{white}{Size} & \textcolor{white}{Description} \\
\midrule
\endhead
Counter & 4 Bytes & Incremental counter to ensure packet freshness. \\
Leader ID & 2 Bytes & ID of the network leader. \\
Counter HMAC & 16 Bytes & HMAC of Counter generated by the leader node. \\
Leader ID HMAC & 16 Bytes & HMAC of Leader ID generated by the leader node. \\
\bottomrule
\end{longtable}

\subsection{Dynamic Connection Tracking}

Static neighbor tables, once populated at deployment or through the initial key-derivation exchange (Section~2.3), do not account for node mobility, intermittent range, or nodes powering on after their neighbours have already converged on a leader. HAVEN addresses this with a lightweight, always-on connection-tracking mechanism layered on top of the existing Neighbor Table (Section~3.3).

\begin{itemize}
\item \textbf{Local Connection Awareness} --- every node continuously maintains a live record of the nodes it is currently in contact with, distinct from the longer-lived, authenticated Neighbor Table entries.
\item \textbf{Ambient Heartbeat Broadcast} --- each node periodically emits a heartbeat that any node within radio range can send or receive, regardless of whether a session already exists between them. Unlike the per-link Leader Heartbeat (Section~3.6), this beacon is not addressed to a specific recipient, but it is still HMAC-authenticated and nonce-protected against replay, per the Protocol Convention (Section~1.4).
\item \textbf{Network Self-Identification} --- the heartbeat carries only the broadcasting node's own Node ID and the Leader ID of the network it currently belongs to. This lets any listener classify the sender at a glance, without a handshake.
\item \textbf{Cross-Network Filtering} --- a node that hears a heartbeat carrying a Leader ID different from its own recognizes the sender as belonging to a foreign network and disregards it for connection-tracking purposes. This prevents nodes from unrelated deployments or adjacent partitions from being mistaken for legitimate neighbours.
\item \textbf{Capacity-Gated Dynamic Admission} --- if a node has spare capacity in its connection table and hears a \emph{verified} heartbeat from an unrecognized node advertising the \emph{same} Leader ID, it dynamically adds that node to its known-neighbour list. Because shared Leader ID implies the two nodes already derive from the same $K_{master}$ (Section~2.3), this admission is a topology update, not a fresh trust decision.
\end{itemize}

\subsubsection*{Presence Heartbeat Packet (24 Bytes)}

\begin{longtable}{@{}>{\bfseries}p{3.4cm}p{2cm}p{8.5cm}@{}}
\toprule
\rowcolor{havennavy}
\textcolor{white}{Field} & \textcolor{white}{Size} & \textcolor{white}{Description} \\
\midrule
\endhead
Node ID & 2 Bytes & ID of the broadcasting node. \\
Leader ID & 2 Bytes & ID of the leader of the broadcasting node's current network. \\
Nonce & 4 Bytes & Freshly generated per broadcast; prevents a captured heartbeat from being replayed to trigger a stale admission. \\
HMAC & 16 Bytes & HMAC over Node ID $\|$ Leader ID $\|$ Nonce, keyed with the broadcasting node's $K_{node}$. \\
\bottomrule
\end{longtable}

\begin{notebox}{Replay Protection}
Because every node in the network already shares $K_{master}$ (Section~2.3), any receiver can independently re-derive the broadcaster's $K_{node} = \text{hmac}(K_{master}, node\_id, 16)$ and verify the HMAC locally, with no prior session or handshake required. The per-broadcast Nonce acts as a monotonically increasing counter, ensuring a captured Presence Heartbeat cannot be recorded and replayed later to falsely re-admit a node that has since moved out of range or been captured; receivers discard any heartbeat whose Nonce is not strictly greater than the last one seen from that Node ID. A heartbeat that fails HMAC verification is discarded outright and never reaches the admission-decision step in Capacity-Gated Dynamic Admission.
\end{notebox}

% ==================== 4. EMBEDDED SYSTEMS ====================
\section{Embedded Systems and Firmware}

\subsection{Microcontroller Architecture}

The HAVEN project targets ARM Cortex-M microcontrollers (specifically the STM32 family, such as Cortex-M0 to M4).\citeref{4} These are single-chip devices containing the CPU, Flash memory, SRAM, and peripherals (SPI, UART, Timers, GPIO) on one die.\citeref{4} There is no operating system; the firmware has complete control.\citeref{4}

\subsubsection*{Memory Layout (32-bit Address Space)}

\begin{longtable}{@{}>{\ttfamily\bfseries}p{3cm}p{10cm}@{}}
\toprule
\rowcolor{havennavy}
\textcolor{white}{\normalfont\bfseries Address} & \textcolor{white}{Region} \\
\midrule
\endhead
0x00000000 & Flash memory (code/.text, .rodata, vector table). \\
0x20000000 & SRAM (volatile memory, .data, .bss, stack). \\
0x40000000 & Peripherals. \\
\bottomrule
\end{longtable}

Peripherals are accessed using memory-mapped I/O, meaning interacting with hardware like the SPI radio chip involves reading and writing to specific memory addresses (e.g., SPI1 registers at \texttt{0x40013008}).\citeref{4} The \texttt{volatile} keyword is crucial in C when accessing these registers to prevent compiler optimization.\citeref{4}

\subsection{Interrupts and Vector Table}

Interrupts allow hardware to pause the CPU's main loop to handle immediate events. The CPU saves its state, runs the Interrupt Service Routine (ISR), and then resumes.\citeref{4} The Vector Table, located at \texttt{0x00000000}, contains the addresses of these handlers.\citeref{4} For example, \texttt{SPI1\_IRQ} at \texttt{0x0800\_0265} handles radio reception, and \texttt{TIM2\_IRQ} at \texttt{0x0800\_0281} handles the election tick timer.\citeref{4}

\subsection{Cross-Compilation and Toolchain}

Firmware is built using cross-compilation on a host PC targeting the ARM architecture.\citeref{4} The toolchain used is \texttt{arm-none-eabi-gcc}.\citeref{4}

\begin{itemize}
\item \texttt{arm} --- target architecture.
\item \texttt{none} --- bare-metal (no OS).
\item \texttt{eabi} --- Embedded Application Binary Interface.
\end{itemize}

Key compiler flags include \texttt{-mcpu=cortex-m4}, \texttt{-Os} (optimize for size), and \texttt{-{}-specs=nosys.specs} (no OS system calls).\citeref{4} A linker script (\texttt{.ld}) explicitly maps code and variables to the correct physical memory addresses.\citeref{4}

\subsection{Hardware Emulation and Debug Communication}

For development and protocol validation, the project has shifted from QEMU to Renode, a multi-node emulation framework capable of instantiating many HAVEN nodes and their radio interconnects within a single simulated test harness --- a closer match to the multi-node gossip and election behaviour described in Section~3 than a single-instance emulator.

Debug and diagnostic output is now handled exclusively through the UART peripheral, across both emulated and physical deployments. This retires the project's earlier use of ARM Semihosting entirely: rather than executing a Breakpoint (\texttt{BKPT}) instruction that a debugger or emulator must intercept, nodes now write debug text directly to a UART transmit buffer, which is read identically whether the target is running under Renode or on physical Cortex-M hardware.

\begin{notebox}{Why UART Only}
Semihosting relied on an attached debugger or emulator to trap the \texttt{BKPT} instruction, and produced a fatal HardFault if that trap was absent --- meaning debug behaviour differed between simulated and physical nodes and could crash a standalone device outright. Standardizing on UART removes this asymmetry: the same debug path works unmodified across Renode simulation and deployed hardware, with no risk to autonomous field nodes.
\end{notebox}

\subsection{Persisting the Epoch Counter}

The \texttt{epoch\_counter} must survive power loss to prevent replay attacks or node lockouts.\citeref{4} It is stored in EEPROM (or emulated EEPROM in flash).\citeref{4} To preserve flash endurance, the counter is only written when the epoch advances, not on every message.\citeref{4}

% ==================== 5. CRYPTOGRAPHIC IMPLEMENTATIONS ====================
\section{Cryptographic Implementations}

\subsection{Python Libraries}

The Python simulations use specific libraries for cryptographic operations:\citeref{4}

\begin{itemize}
\item \texttt{hashlib.sha256()} --- core primitive for integrity checks and HMACs.
\item \texttt{hmac.new()} --- binds data to a secret key.
\item \texttt{hmac.compare\_digest()} --- prevents timing attacks by comparing tags in constant time.
\item \texttt{secrets.token\_bytes()} and \texttt{secrets.randbelow()} --- generate cryptographically strong random values for keys, nonces, and DH exchanges.
\item \texttt{struct.pack()} / \texttt{struct.unpack()} --- packs integers into binary buffers.
\end{itemize}

\subsection{mbedTLS on Cortex-M}

The physical firmware utilizes the mbedTLS library.\citeref{4} The \texttt{haven\_hmac()} function wraps \linebreak\texttt{mbedtls\_md\_hmac} to generate various keys and signatures:

\begin{itemize}
\item \textbf{$K_{node}$ derivation} --- \texttt{hmac(K\_master, node\_id, 16)}.
\item \textbf{Priority score} --- \texttt{hmac(K\_node, node\_id XOR epoch, 16)}.
\item \textbf{HELLO MAC} --- \texttt{hmac(K\_AB, HELLO\_RESP fields, 16)}.
\item \textbf{CLAIM link\_mac} --- \texttt{hmac(K\_link[j], all preceding fields, 16)}.
\end{itemize}

% ==================== 6. TECHNOLOGY STACK ====================
\section{Technology Stack}

\begin{longtable}{@{}>{\bfseries}p{3.4cm}p{10.2cm}@{}}
\toprule
\rowcolor{havennavy}
\textcolor{white}{Technology} & \textcolor{white}{Role} \\
\midrule
\endhead
Python & Simulating network protocols, fine-tuning parameters, and network modeling using the NetworkX library (e.g., Random Geometric Graphs). \\
Bare-metal C & Firmware programming. \\
Renode & Multi-node hardware emulation, including radio-link simulation across many HAVEN nodes in a single test harness. \\
mbedTLS & Submodule for HMAC operations on firmware. \\
CMake & Build system automation for cross-compilation. \\
NS-3 & Planned for large-scale network simulation. \\
\bottomrule
\end{longtable}

% ==================== 7. GIT REPOSITORY REFERENCE ====================
\section{Git Repository Reference}

The project codebase is managed via Git.\citeref{2} The repository includes branches such as \texttt{main}, \texttt{dev}, \texttt{dev/test}, and specific feature branches like \texttt{turboblom-python}, \texttt{turboblom-c}, \texttt{qemu-bringup}, and \texttt{node-capture}.\citeref{2}

The codebase is structured with directories for simulation (\texttt{sim/}) containing Python scripts (e.g., \texttt{turboblom.py}, \texttt{network\_sim.py}) and firmware (\texttt{firmware/}) containing C source files, linker scripts, and CMake configurations.\citeref{2}

\end{document}
